\section{Feasible bases}\label{sec:feasibility}

We now apply the holomorphic mass estimate to finitely many real strip
constraints. It gives a lower bound for the sum of their boundary
feasibility probabilities. For the lens of Lemma~\ref{lem:planar-lens},
these probabilities will become volumes of real simplices.

Fix $d\geq1$ and nonzero real row vectors $b_1,\ldots,b_m$ spanning
$(\R^d)^*$. Write
\begin{equation}\label{eq:strip-body}
 A=\{X\in\R^d:|b_iX|\leq1\text{ for every }i\},
 \qquad C=A^\circ=\operatorname{conv}\{\pm b_1,\ldots,\pm b_m\}.
\end{equation}
We identify rows with vectors when taking a convex hull, and extend their
action complex-linearly to $\C^d$. The spanning assumption makes $A$
bounded, and $0\in\operatorname{int}A$. The polar identity follows from
the bipolar theorem: the polar of the displayed convex hull is exactly
$A$.
Let $\mathcal B$ be the collection of $d$-element sets of independent rows.
For $I\in\mathcal B$, put $B_I=(b_i)_{i\in I}$ in increasing index order
and $D_I=|\det B_I|$. Dependent row sets are omitted throughout.

For the next lemma, $F$ may be any biholomorphism of the unit disk onto a
bounded planar domain $D$, with $F(0)=0$ and a continuous extension to the
closed disk. For independent uniform angles $\theta_i\in[0,2\pi)$,
$i\in I$, define
\begin{equation}\label{eq:feasible-probabilities}
 Z_I(\theta)=B_I^{-1}\bigl(F(e^{i\theta_i})\bigr)_{i\in I},
 \qquad
 P_I=\mathbb P\{b_jZ_I(\theta)\in\overline D\text{ for all }j\},
 \qquad S=\sum_{I\in\mathcal B}P_I.
\end{equation}
The active coordinates $i\in I$ already belong to $\overline D$; thus the
probability tests only the remaining rows.

\begin{lemma}[Boundary feasibility]\label{lem:holomorphic-feasibility}
Let $d\ge1$ and let $b_1,\ldots,b_m$ be nonzero real rows spanning
$(\R^d)^*$. Let $F:\mathbb D\to D$ be a biholomorphism onto a bounded
planar domain, continuous on $\overline{\mathbb D}$, with $F(0)=0$.
Then the sum $S$ in \eqref{eq:feasible-probabilities} satisfies $S\ge1$.
No condition on extra-row boundary equalities or on dependent row
subsets is required.
\end{lemma}

\begin{proof}
Let $g=F^{-1}$ and set
\[
 \Omega=\{z\in\C^d:b_jz\in D\text{ for every }j\},
 \qquad h_j(z)=g(b_jz).
\]
The set $\Omega$ is open, bounded, and contains $0$. For an integer
$k\geq1$, apply Lemma~\ref{lem:holomorphic-mass} to the holomorphic map
\[
 f_k=(h_1^k,\ldots,h_m^k),\qquad
 \tau_k=\sum_{j=1}^m|h_j|^{2k}.
\]
We shall let $k$ tend to infinity. The mass lower bound has scale
$k^d$. The radial substitution below removes the same factor from the
Jacobian integrals and leaves a fixed domain on which the inferred
points tend to boundary samples.
The unique zero of $f_k$ is $0$: $g$ vanishes only at $0$, and the rows span.
Writing $c=g'(0)\ne0$, its degree-$k$ leading term is
$((cb_1z)^k,\ldots,(cb_mz)^k)$, also nonzero whenever $z\ne0$.
For $0<s<1$, the set $\{\tau_k\leq s\}$ is compact in $\Omega$.
Indeed each of its row coordinates belongs to the fixed compact set
$F(\{|w|\leq s^{1/(2k)}\})\subset D$. A basis of rows bounds $z$;
limits retain all these inclusions and the inequality $\tau_k\leq s$.
All hypotheses of the mass lemma therefore hold.
The derivative row of $h_j^k$ is a scalar multiple of $b_j$, so dependent
row subsets have zero minors. Consequently
\begin{equation}\label{eq:feasibility-mass}
 \sum_{I\in\mathcal B}M_{I,k}\geq
 k^d\frac{\pi^d}{d!},\qquad
 M_{I,k}:=\int_{\{\tau_k<1\}}
 \left|\det\frac{\partial(h_i^k)_{i\in I}}
                        {\partial(z_1,\ldots,z_d)}\right|^2\,dV_{2d}(z).
\end{equation}

Fix $I$. The change of variables $u_i=h_i(z)$, $i\in I$, is one-to-one
onto its image, with inverse
$z=B_I^{-1}(F(u_i))_{i\in I}$. Its squared complex Jacobian is its real
volume Jacobian. The squared determinant in $M_{I,k}$ is this Jacobian factor
times $k^{2d}\prod_i|u_i|^{2k-2}$.
The transformed integration domain is contained in
\[
 \left\{u\in\mathbb D^d:
       \sum_{i\in I}|u_i|^{2k}<1,\quad
       b_jB_I^{-1}(F(u_i))_{i\in I}\in D\text{ for all }j\right\}.
\]
We use this larger domain for an upper bound, and then put
$u_i=\rho_i^{1/k}e^{i\theta_i}$. The coordinate hyperplanes, on which
polar coordinates are undefined, have Lebesgue measure zero. Since
\[
 k^2|u_i|^{2k-2}\,dV_2(u_i)
      =k\rho_i\,d\rho_i\,d\theta_i,
\]
we obtain
\begin{equation}\label{eq:feasibility-radial}
 \frac{M_{I,k}}{k^d}\leq
 \int_{\substack{\rho_i>0,\ \sum_i\rho_i^2<1\\
                  \theta\in[0,2\pi)^d}}
 \mathbf1\!\left\{
  b_jB_I^{-1}\bigl(F(\rho_i^{1/k}e^{i\theta_i})\bigr)_{i\in I}
                    \in D\text{ for every }j\right\}
 \prod_{i\in I}\rho_i\,d\rho_i\,d\theta_i.
\end{equation}

The measure in this integral is finite: its total mass is the volume
$\pi^d/d!$ of the unit ball in $\C^d$. For every positive $\rho_i$,
continuity of $F$ on the closed disk makes the inferred points converge
to $Z_I(\theta)$. If a limiting row coordinate lies outside
$\overline D$, that constraint fails for all sufficiently large $k$.
Thus the upper limit of the indicator is bounded by the indicator of
the event defining $P_I$, including all possible boundary equalities.
The upper-bound form of Fatou's lemma, applied to functions bounded by
$1$ on this finite measure space, gives
\[
 \limsup_{k\to\infty}\frac{M_{I,k}}{k^d}
 \leq \frac{\pi^d}{d!}P_I.
\]
There are only finitely many bases, so \eqref{eq:feasibility-mass}
and these upper bounds give
\begin{align*}
 \frac{\pi^d}{d!}
 &\leq\limsup_{k\to\infty}
          \sum_{I\in\mathcal B}\frac{M_{I,k}}{k^d}\\
 &\leq\sum_{I\in\mathcal B}
          \limsup_{k\to\infty}\frac{M_{I,k}}{k^d}\\
 &\leq\frac{\pi^d}{d!}\sum_{I\in\mathcal B}P_I.
\end{align*}
Thus $S\geq1$.
\end{proof}

\subsection{Boundary equalities for the lens}

The lower bound for $S$ allows additional boundary equalities. The real
volume identity in Section~\ref{sec:simplex} will need the stronger fact
that these equalities have probability zero.
From now on, $F,D$, and $\lambda$ are those of
Lemma~\ref{lem:planar-lens}. We may also require that no two rows are
proportional. For any finite strip presentation this is achieved without
changing $A$ by keeping, on each row line, a row of largest magnitude;
its strip implies all the discarded strips on that line. Zero rows, if
present initially, impose no constraint and are discarded.
In dimension one the reduced list has a single row, so there are no
extra-row boundary equalities to consider.

\begin{lemma}[Additional boundary equalities have probability zero]
\label{lem:boundary-ties}
Let $F:\mathbb D\to D$ be the lens map of Lemma~\ref{lem:planar-lens}.
Assume that the nonzero spanning rows $b_1,\ldots,b_m$ are pairwise
nonproportional. For every $I\in\mathcal B$ and $j\notin I$,
\[
 \mathbb P\{b_jZ_I(\theta)\in\partial D\}=0.
\]
\end{lemma}

\begin{proof}
Write $b_jB_I^{-1}=(c_i)_{i\in I}$. At least two coefficients, say
$c_a,c_b$, are nonzero; otherwise $b_j$ would be proportional to a basis
row. Let $\gamma(\theta)=F(e^{i\theta})$.
Apart from its two endpoints, the lens boundary consists of the smooth
regular branches $(\pm\lambda(t),t)$, $-1<t<1$.
Their tangent directions have at most one occurrence on each branch for
any fixed unoriented line direction, because $\lambda'$ is strictly
monotone. The boundary parameter has constant nonzero speed in $t$ on
each branch, as established in the proof of Lemma~\ref{lem:planar-lens}.

Hold all phases other than $\theta_a,\theta_b$ fixed. The map
\[
 (\theta_a,\theta_b)\longmapsto
   c_a\gamma(\theta_a)+c_b\gamma(\theta_b)
       +\sum_{i\ne a,b}c_i\gamma(\theta_i)
\]
has real Jacobian determinant
$\det(c_a\gamma'(\theta_a),c_b\gamma'(\theta_b))$.
For every nonendpoint $\theta_a$, it vanishes at only finitely many
nonendpoint values of $\theta_b$. Its zero set, including the excluded
endpoint lines, therefore has planar measure zero by Fubini's theorem.
Off that set the inverse function theorem gives local diffeomorphism
charts. A countable subcover of such charts suffices, and in each chart
the inverse is locally Lipschitz, so it sends a planar null set to a null
set. The set $\partial D$ is planar-null, being a countable union of
compact smooth subarcs and two points. Its inverse image consequently
has measure zero in the two remaining phases. A final application of
Fubini over the fixed phases proves the assertion.
\end{proof}
